%% ---- 2. Trabalhos Relacionados ---------------------------------------------
\section{Trabalhos Relacionados}

\paragraph{Explicadores com acesso ao modelo.}
Em recomendação, o ACCENT~\cite{tran2021accent} usa funções de influência, o
PRINCE~\cite{ghazimatin2020prince} explora o grafo de interações e o
CountER~\cite{tan2021counter} otimiza perturbações contínuas com gradiente. Em
visão, DiME~\cite{jeanneret2022dime} e LDCE~\cite{schrodi2024ldce} guiam um
modelo de difusão pelo gradiente do classificador. Todos exigem acesso interno
ao modelo ou entrada contínua; aqui o histórico é discreto e o recomendador só
devolve escores.

\paragraph{Caixa-preta em recomendação sequencial.}
O GECE~\cite{scarcelli2025gece} é o trabalho mais próximo: um algoritmo
genético sobre a sequência de itens, com mutações de substituição, inserção e
remoção, cuja aptidão é a soma ponderada da distância de edição e de uma
diferença de escores. É monoobjetivo e devolve um contrafactual. O
TRACE~\cite{kim2026trace} gera contrafactuais direcionados por busca em feixe
com orçamento de consultas e traz uma implementação pública do GECE, que usamos
como linha de base (Seção~\ref{sec:base}). Os dois editam a sequência, isto é,
escolhem também \emph{qual item} colocar em cada posição; nós buscamos apenas
\emph{quais posições} importam e sorteamos os itens substitutos.

A Tabela~\ref{tab:related} resume a posição deste trabalho. O CountER, com
acesso ao gradiente, e o GECE, de caixa-preta, agregam os critérios em um
escalar e devolvem uma única resposta; este trabalho mantém o acesso apenas a
escores, como o GECE, mas trata os critérios como um vetor e devolve uma
frente de Pareto.

\paragraph{Necessidade, suficiência e frentes de Pareto.}
As probabilidades de necessidade e de suficiência vêm de
Pearl~\cite{pearl2009causality} e foram usadas em recomendação pelo
CountER~\cite{tan2021counter}, que as combina em uma otimização escalar. O
MOC~\cite{dandl2020moc} aplica o NSGA-II a contrafactuais em dados tabulares e
mostra que a frente de Pareto é uma saída útil. Não encontramos equivalente
sobre históricos de itens, mas essa é uma afirmação negativa fraca.
Mohammadi et al.~\cite{mohammadi2025beyondtop1} mostram que a avaliação de
explicações contrafactuais em recomendação é inconsistente entre trabalhos, o
que motivou o cuidado com orçamento igual e fundos comuns a todos os métodos.

\begin{table}[t]
\caption{Posição deste trabalho. ``Escalar'' = critérios somados em um único valor.}
\label{tab:related}
\small
\begin{tabularx}{\columnwidth}{@{}lXXX@{}}
\toprule
 & \textbf{CountER} & \textbf{GECE} & \textbf{Este trabalho} \\
\midrule
Acesso      & gradiente & escores & escores \\
Busca       & otimização contínua & AG & NSGA-II \\
Critérios   & escalar & escalar & vetor $(N,S,C)$ \\
Saída       & 1 explicação & 1 sequência & frente de Pareto \\
\bottomrule
\end{tabularx}
\end{table}
